\section {Условная вероятность. Независимость событий. Критерий независимости.}

\begin{defn}
	Пусть задано вероятностное пространство $(\Omega, \SigAlg, \MyPr), \\ \text{события } {A, B \in \SigAlg,~ \MyPr(B) > 0}$. 
	\textit{Условная вероятность события $A$ при событии $B$}:
	\begin{equation*}
        \MyPr(A|B)=\frac{\MyPr(AB)}{\MyPr(B)}.
	\end{equation*}
\end{defn}

\begin{thm*}
	Условная вероятность $\MyPr(A|B)$~--- вероятность, заданная на $\sigma$-алгебре $\SigAlg$.
\end{thm*}

\begin{proof}
	Проверим три аксиомы из определения вероятности.
	
	\begin{enumerate}
		\item 
		      $\forall A \in \SigAlg \;\; \MyPr(A|B) \geqslant 0, \text{т.к.}~ \MyPr(AB) \geqslant 0,~ \MyPr(B) > 0$;
		\item 
		      $\MyPr(\Omega|B) = \cfrac{\MyPr(B \cap \Omega)}{\MyPr(B)} = \cfrac{\MyPr(B)}{\MyPr(B)} = 1$;
		\item 
              Пусть дана некоторая последовательность событий $A_1, A_2, \ldots, A_n, \ldots$; $A_i A_j = \varnothing~ (i \ne j)$. 
		      Тогда: 
		      \begin{multline*}
			      \MyPr\left(\,\left. \bigcup\limits_{i=1}^\infty A_i \right| B\right)
                  = \frac{\MyPr\left(\left(\, \bigcup\limits_{i=1}^\infty A_i\right) \cap B \right)}{\MyPr(B)}
                  = \frac{\MyPr\left(\, \bigcup\limits_{i=1}^\infty\left(A_i \cap B\right) \right)}{\MyPr(B)} = \\
                  = \frac{\sum\limits_{i=1}^\infty \MyPr(A_i \cap B)}{\MyPr(B)}
			      = \sum\limits_{i=1}^\infty \MyPr(A_i | B). \qedhere
		      \end{multline*}
	\end{enumerate}
\end{proof}
\begin{rmrk}
	Некоторые свойства условной вероятности:
	\begin{enumerate}
		\item 
		      Если $A \cap B = \varnothing,$ то $\MyPr(A | B) = 0.$ 
		\item 
		      Если $B \subset A$, то $\MyPr(A|B) = 1.$ 
		      Например, $\MyPr(B|B) = 1.$
	\end{enumerate}
\end{rmrk}

\subsubsection{Независимость событий}
\begin{defn}
	Пусть есть вероятностное пространство $(\Omega, \SigAlg, \MyPr)$. 
	События $A_1, \ldots, A_n \in \SigAlg$ называются \textit{независимыми в совокупности}, если 
	$\forall k = \overline{2, n}$ $\forall i_{1}, \ldots, i_{k} \colon 1 \leqslant i_1 < i_2 < \ldots < i_k \leqslant n$ выполняется:
	\begin{equation*}
		\MyPr\left(\, \bigcap\limits_{j=1}^k A_{i_j}\right) = \prod\limits_{j=1}^k \MyPr(A_{i_j}).
	\end{equation*}
	
	Иными словами, события независимы в совокупности, если вероятность одновременного наступления \textit{любого набора} из этих событий равна \textit{произведению вероятностей} событий, входящих в этот набор. 
	В частности, при $n = 2$ события $A$ и $B$ независимы, если $\MyPr(AB) = \MyPr(A)\MyPr(B)$.
\end{defn}

\begin{namedthm}[Свойства независимых событий]\leavevmode
	\begin{enumerate}
		\item 
		      Если $A = \varnothing$ или $\MyPr(A) = 0$, то $\forall B \colon \MyPr(B) > 0$ события $A$ и $B$ независимы.
		\item 
		      Пусть $A$ и $B$ независимы. Тогда события $\overline{A}$ и $B$, $A$ и $\overline{B}$, $\overline{A}$ и $\overline{B}$ также независимы. 
		\item 
		      Пусть $A \subset B$ и $\MyPr(A) > 0, \, \MyPr(B) < 1$. Тогда $A$ и $B$ зависимы. 
		\item 
		      Если события $A$ и $B$ независимы и $\MyPr(B) > 0$, то $\MyPr(A|B) = \MyPr(A)$.
	\end{enumerate}
\end{namedthm}

\begin{proof}
	\begin{enumerate} 
		\item 
		      Если $A = \varnothing$, то $AB = \varnothing \implies \MyPr(AB) = 0.$ Но $ \MyPr(A)\MyPr(B) = 0 \cdot \MyPr(B) = 0 \implies \MyPr(AB) = \MyPr(A) \MyPr(B)$.
		              
		      Если же ${\MyPr(A) = 0}$, то ${AB \subset A \implies \MyPr(AB) \leqslant \MyPr(A) = 0}$. 
		      В то же время ${0 = \MyPr(AB) = 0 \cdot \MyPr(B) = \MyPr(A) \MyPr(B)}$.
		\item 
		      Докажем независимость $\overline{A}$ и $B$, представив последнее в виде \\
		      $B = AB \cup \overline{A}B$. Тогда
		      \begin{multline*}
		      	\MyPr(B) = \MyPr(AB) + \MyPr(\overline{A}B) = \MyPr(A)\MyPr(B) + \MyPr(\overline{A}B) \implies \\
                  \implies \MyPr(\overline{A}B) = \MyPr(B) - \MyPr(A)\MyPr(B) = \MyPr(B) (1 - \MyPr(A)) = \MyPr(\overline{A})\MyPr(B).
		      \end{multline*}
		      Независимость $\overline{A}$ и $B$ доказана. 
		      Аналогично доказываются остальные утверждения.
		\item
		      Предположим, что события независимы.
		      Тогда $\MyPr(AB) = \MyPr(A)\MyPr(B),$ но в силу вложенности $A \subset B$: $\MyPr(AB) = \MyPr(A) > 0$, следовательно, $\MyPr(B) = 1$, что противоречит условию.
		\item
		      $\MyPr(A | B) = \cfrac{\MyPr(AB)}{\MyPr(B)} = \cfrac{\MyPr(A)\MyPr(B)}{\MyPr(B)} = \MyPr(A)$.
		      \qedhere
	\end{enumerate}
\end{proof}

\begin{rmrk}
	В общем случае из попарной независимости событий $A_1, \ldots, A_n$ не следует их независимость в совокупности.
	\begin{exmp}
		Рассмотрим правильный тетраэдр, три грани которого окрашены соответственно в красный, синий, зелёный цвета, а четвёртая грань содержит все три цвета. 
		Событие $R$ (соответственно, $G$, $B$) означает, что выпала грань, содержащая красный (соответственно, зелёный, синий) цвета.
		
		Т.к. каждый цвет есть на двух гранях из четырёх, то
		\begin{equation*}
            \MyPr(R) = \MyPr(G) = \MyPr(B) = \frac{1}{2}.
		\end{equation*}
		Вероятность пересечения, соответственно:
		\begin{equation*}
            \MyPr(RG) = \MyPr(GB) = \MyPr(RB) = \frac{1}{4} = \frac{1}{2} \cdot \frac{1}{2},
		\end{equation*}
		следовательно, все события попарно независимы. 
		Однако вероятность пересечения всех трёх:
		\begin{equation*}
            \MyPr(RGB)  = \frac{1}{4} \neq \MyPr(R) \MyPr(G) \MyPr(B),
		\end{equation*}
		т.е. события не являются независимыми в совокупности. 
	\end{exmp}
\end{rmrk}

\begin{symb}
	\begin{equation*}
		A_{i}^{(\delta)} =
		\begin{cases}
			A_{i},            & \delta = 1; \\
			\overline{A_{i}}, & \delta = 0. 
		\end{cases}
	\end{equation*}
\end{symb}

\begin{namedthm}[Критерий независимости]
    События $A_1, \ldots, A_n$ независимы в совокупности $\iff \forall \delta_1, \delta_2, \ldots, \delta_n \in \{0, 1\}$ выполнено равенство
	\begin{equation*}
		\MyPr\left(\, \bigcap_{i=1}^{n} A_{i}^{\left( \delta_{i} \right)} \right)
		= \prod_{i=1}^{n}\MyPr\left( A_{i}^{\left(\delta_{i}\right)} \right).
	\end{equation*}
\end{namedthm}
